% Nouveau EGA V (cote 25): archive cote 25, site page 93 (PDF page 94).
% typescript. Draft compared: none.
% AI-transcribed by Claude Opus 5.5 (claude-opus-5-5, adaptive thinking, effort high), Message Batches API, batch ega25b (msgbatch_01A7n5YuXubUPNK1zsepi8RG), 2026-11-13; not yet checked by eye.
\scanpage{25/93}
\mspage{87}
de G, on peut également prendre un point à valeurs dans un \inserted[typed]{S-}préschéma T,
on obtient alors $\underline{P}_{\hand{\unsure{\xi|\zeta}}} \longrightarrow \unsure{P}_T$ (T-morphisme, qui est une immersion ferm\unsure{ée|é}
dans le cas envisagé d'abord).

Si $f:X \longrightarrow P$ est un morphisme, on pose $\underline{X} = X\times_P\underline{P}$ , et on obtient
un diagramme de même type que le carré précédent
\diagram{\begin{tikzcd}
X \arrow[d] & \underline{X} \arrow[l, dash] \arrow[d] \\
S & G \arrow[l]
\end{tikzcd}}
.

Il est alors évident que toutes les questions étudiées dans les n$^{\text{os}}$ 2 à 8
gardent un sens dans le contexte général qu'on vient de poser, et il y a
lieu de dégager les conditions axiomatiques qui assurent les conclusions
tirées dans lesdits n$^{\text{os}}$ .

Nous allons supposer $\underline{P}$ \uline{et} G \uline{plats sur} S, G étant \uline{à fibres}
\inserted[typed]{\xout{de dimen}}
\uline{géométriquement irréductibles} (pour pourvoir en considérer les points
génériques !), \uline{de dimension} N, \uline{le morphisme} $\underline{P} \longrightarrow P$ \uline{est supposé lisse}
\uline{à fibres géométrique\xout{x}ment irréductibles de dimension N-m} . Donc le
morphisme $\underline{X} \longrightarrow X$ aura les mêmes propriétés. Toutes les propriétés
envisagées \inserted{ici et plus tard,} sont stables par changment de base sur S, et s'applique\xout{x}ront
en particulier aux fibres.

Supposons d'abord S=Spec(k). Soit Z une partie fermée de X de
dimension d, alors son image inverse $\underline{Z}$ dans $\underline{X}$ est une partie fermée de
dimension $d + (N-m) = N + d-m$ . Si $d < m$, $\underline{Z}$ est de dimension $< N$,
donc $\underline{Z} \longrightarrow G$ ne peut être dominant, donc si \hand{$\eta$} est le point générique de
G, on a $\underline{Z}_{\hand{\eta}} = \emptyset$ ; en fait, ce raisonnement montre même (en remplaçant
Z par $\overline{f(Z)}$ ) que si dim f(Z) \hand{$<$} m, alors $\underline{Z}_{\hand{\eta}} = \emptyset$ . On veut une condi-%
tion \inserted{sur (D)} assurant que si dim f(Z) \over{\unsure{$<$}}{$\geq$} m, alors $\underline{Z}_{\hand{\eta}} \neq \emptyset$ ; \deleted{et de façon plus
précise que dim $\underline{Z}$ = dim Z - m . Pour la propriété $\underline{Z} \neq \emptyset$ ,} il

\mspage{93}
